% Part: counterfactuals
% Chapter: minimal-change-semantics
% Section: introduction

\documentclass[../../../include/open-logic-section]{subfiles}

\begin{document}

\olfileid{cnt}{min}{int}

\olsection{引言}

Stalnaker 和 Lewis 对诸如“如果火柴被划着，它就会点燃”这样的反事实条件句提出了阐释。他们的阐释旨在说明如何正确理解这类语句的真值条件。这两种提议背后的核心思想是：要评估一个反事实条件句是否为真，我们必须考虑那些与实际世界差异最小、且使得前件为真的可能世界。如果后件在这些可能世界中为真，那么该反事实条件句就为真。例如，假设我手里拿着一根火柴和一个火柴盒。在实际世界里，我只是看着它们，设想如果我划着火柴会发生什么。从实际世界到那个我划着火柴的世界，所需的最小改变就是我决定行动并划着火柴。它之所以是最小的，在于其他一切都没有改变：我没有同时跳起来，划火柴也不会顺带点燃我的头发，我的手指没有突然失去所有力气，我也没有在遭到水枪伏击的同时被浇湿，等等。在那个替代的可能性中，火柴点燃了。因此，“如果我要划着火柴，它就会点燃”是真的。

这种直观阐释可以与反事实条件句逻辑的形式语义相配合。Lewis 为反事实条件句引入了符号“$\boxright$”，而 Stalnaker 则使用了符号“$>$”。我们将采用~$\cif$，并将其作为一个二元联结词添加到命题逻辑中。因此，除了形如 $!A \lif !B$ 的公式外，我们还有形如~$!A \cif !B$ 的公式。其形式语义类似于模态逻辑的关系语义，基于在世界上对公式求值的模型，而定义 $\mSat{M}{!A \cif !B}[w]$ 的满足条件则是由某些（其他）世界~$w'$ 上的 $\mSat{M}{!A}[w']$ 和 $\mSat{M}{!B}[w']$ 给出的。哪些 $w'$ 呢？直观地说，是那些最接近~$w$ 且使得~$\mSat{M}{!A}[w']$ 成立的世界。这就要求“接近”关系也必须被包含在模型之中。

Lewis 引入了一种用图形表示反事实情境的富有启发性的方法。每个可能世界都位于一组包含其他世界的嵌套球体的中心——我们将这些球体画成同心圆。两个球体之间的世界彼此同样接近中心世界，包含在内层球体中的世界更近，而位于外围球体中的世界更远。

\begin{center}
\begin{tikzpicture}[scale=.7]\small
  \spheresystem{5}
  \draw (0,0) node {$w$};
  \propositionintersect{0}{4}{40}{3.5}
  {\tikzset{proposition/.append={smooth,tension=1.4}}
  \proposition[shift={(1.6,-1)}]{90}{1}{30}{4}}
  \path (1.5,3) node[below] {$\formula{B}$};
  \path (3,0) node {$\formula{A}$};
\end{tikzpicture}
\end{center}

最近的 $!A$-世界是那些满足~$!A$ 并位于中心世界~$w$ 周围最小球体（灰色区域）内的世界。直观上，如果 $!B$ 在所有最近的 $!A$-世界中都为真，则 $!A \cif !B$ 在~$w$ 处得到满足。
\end{document}